\section{Knowledge Graphs and Ontologies}
\label{sec:gr-kg-ontologies}

A knowledge graph represents knowledge as a graph whose nodes are entities of
interest and whose edges denote relationships between them, providing an explicit
and machine-readable model of a domain~\cite{hogan2021kg}. Unlike free text,
this structure makes relations first-class objects that can be queried,
validated, and reasoned over~\cite{hogan2021kg, ji2021kgsurvey}. Knowledge graphs
have been adopted at scale by both industry and academia precisely because they
can integrate diverse, dynamic, large-scale collections of
data~\cite{hogan2021kg}.

\subsection{Entities, Relations, and Triples}

The atomic unit of a knowledge graph is the \emph{triple}: a subject entity, a
predicate (relation), and an object, written as (subject, predicate, object),
which asserts a single labelled edge between two nodes~\cite{hogan2021kg}. A
collection of such triples forms a directed, edge-labelled graph, and this simple
model underlies the surveyed families of graph data
models~\cite{hogan2021kg, angles2008graphdb}. Representing knowledge as triples
is what allows a graph to be traversed to answer relational and multi-hop
questions that flat text cannot support directly~\cite{ji2021kgsurvey}.
\Cref{fig:gr-kg-triples} shows a small example graph assembled from such triples.

\begin{figure}[htbp]
  \centering
  \begin{tikzpicture}[
    ent/.style={draw, ellipse, minimum width=2.4cm, minimum height=1cm,
                align=center, font=\small},
    rel/.style={-{Latex[length=2mm]}, thick, font=\scriptsize}]
    \node[ent] (r1) {DeepSeek-R1};
    \node[ent, right=3.4cm of r1] (rm) {Reasoning\\Model};
    \node[ent, below=2.2cm of r1] (math) {MATH\\Dataset};
    \node[ent, below=2.2cm of rm] (student) {Student\\Model};
    \draw[rel] (r1) -- (rm)      node[midway, above] {\texttt{isA}};
    \draw[rel] (r1) -- (math)    node[midway, left]  {\texttt{trainedOn}};
    \draw[rel] (r1) -- (student) node[midway, above, sloped] {\texttt{teaches}};
    \draw[rel] (student) -- (rm) node[midway, right] {\texttt{isA}};
  \end{tikzpicture}
  \caption[An example knowledge graph of triples]{A knowledge graph is a set of (subject, predicate, object) triples,
    each asserting one labelled edge between two entities~\cite{hogan2021kg}.}
  \label{fig:gr-kg-triples}
\end{figure}

\subsection{Ontology Components: Classes, Properties, and Individuals}

Whereas a knowledge graph records instance-level facts, an \emph{ontology}
specifies the schema those facts must obey. Formally, an ontology is an explicit
specification of a conceptualization---a description of the concepts and
relationships that can exist in a domain~\cite{gruber1993ontology}. Its core
building blocks are classes (types of entity), properties (the relations and
attributes that may hold), and individuals (the concrete instances), together
with axioms that constrain how they may be combined~\cite{hogan2021kg}. By fixing
this vocabulary, an ontology gives the graph a shared, consistent meaning and
enables logical inference over it~\cite{gruber1993ontology, hogan2021kg}.

\subsection{Standards: RDF, RDFS, and OWL}

These ideas are standardized by the World Wide Web Consortium as part of the
Semantic Web vision, in which data is published with well-defined meaning so that
machines can process it~\cite{bernerslee2001semantic}. The Resource Description
Framework (RDF) provides the triple-based data model and abstract syntax for
representing graph data on the Web~\cite{w3c2014rdf}, while the Web Ontology
Language (OWL) supplies the expressive constructs---class hierarchies, property
characteristics, and logical axioms---needed to define rich ontologies and to
support automated reasoning~\cite{w3c2012owl}. Together they form the standard
stack on which interoperable knowledge graphs are built~\cite{hogan2021kg}.
